\begin{figure}[t]
\begin{minted}[fontsize=\small,frame=single]{text}
// OpenQASM 3 teleportation example
OPENQASM 3.0;
include "stdgates.inc";

// Declare qubits and classical bits
qubit[3] q;
bit[2] c;

// Create entanglement between q[0] (alice) and q[1] (bob)
h q[0];
cx q[0], q[1];

// Alice's Bell measurement with msg qubit q[2]
cx q[2], q[0];
h q[2];

// Measure and store results
c[0] = measure q[2];
c[1] = measure q[0];

// Real-time classical control: apply corrections to bob
if (c[1] == 1) x q[1];
if (c[0] == 1) z q[1];
\end{minted}
\caption{Quantum teleportation in OpenQASM 3. Variables are explicitly declared with storage classes (\texttt{qubit[3]}, \texttt{bit[2]}). The \texttt{if} statements demonstrate real-time classical control: corrections are applied to the target qubit based on measurement outcomes within the same execution context. Gates from \texttt{stdgates.inc} (\texttt{h}, \texttt{cx}, \texttt{x}, \texttt{z}) operate on individual array elements. This low-level representation serves as a common compilation target for higher-level languages.}
\label{fig:openqasm-teleport}
\end{figure}
